\input{../header}

\begin{document}
%


%\onehalfspacing
\allowdisplaybreaks
%##################################################################
\section{PS\#14 - The normal distribution! -- {\color{red} KEY} }

\everymath{\displaystyle}
\begin{enumerate}[leftmargin=0pt]
    \item Draw a graph of $p(x)$. 
    \item (Shade in the region of integration.)
    
    \item (Determine bounds of integration.) \begin{red} From $a=-1$ to $b=1$.
    \end{red}
    \item (Not integrable.)
    \item Write down the Taylor series for $f(u) = e^u$.
    
    \begin{red}
    \[e^u = \sum_{k=0}^\infty \frac{1}{k!} u^k = 1 + \frac{u}{1!} + \frac{u^2}{2!} + \frac{u^3}{3!} + \ldots \]
    \end{red}
    \item Use the Taylor series for $e^u$ to write down the Taylor series for $p(x).$
    \begin{red}
    \[T_p(x) = \frac{1}{\sqrt{2\pi}} \sum_{k=0}^\infty \frac{(-1)^k x^{2k}}{k!\cdot 2^k} = \frac{1}{\sqrt{2\pi}} \left(1 - \frac{x^2}{1! \cdot 2} + \frac{x^4}{2!\cdot 4} - \frac{x^6}{3!\cdot 8} + \frac{x^8}{4! \cdot 16} \ldots \right) \]
    \end{red}
    \item Use the Taylor series for $p(x)$ to write down a Taylor series for its antiderivative $P(x)$.    
    \begin{red}    
    \[T_P(x) = \frac{1}{\sqrt{2\pi}} \sum_{k=0}^\infty \frac{(-1)^k x^{2k+1} }{k!\cdot 2^k \cdot (2k+1)} = \frac{1}{\sqrt{2\pi}} \left(x - \frac{x^3}{1! \cdot 2\cdot 3} + \frac{x^5}{2!\cdot 4\cdot 5} - \frac{x^7}{3!\cdot 8 \cdot 7} + \frac{x^9}{4! \cdot 16\cdot 9} \ldots \right)\]
    \end{red}
    \item Use the first four nonzero terms of the Taylor series for $P(x)$ to approximate $P(b)$.     
    \begin{red}
    \[T_P(1) = \frac{1}{\sqrt{2\pi}} \left(1 - \frac{1^3}{1! \cdot 2\cdot 3} + \frac{1^5}{2!\cdot 4\cdot 5} - \frac{1^7}{3!\cdot 8 \cdot 7}\right) = \frac{1}{\sqrt{2\pi}}\cdot  \frac{479}{560} \approx 0.34124.\]
    \end{red}
    \item Use the first four nonzero terms of the Taylor series for $P(x)$ to approximate $P(a)$.    
    \begin{red}
    \[T_P(-1) = \frac{1}{\sqrt{2\pi}} \left(-1 + \frac{1^3}{1! \cdot 2\cdot 3} - \frac{1^5}{2!\cdot 4\cdot 5} + \frac{1^7}{3!\cdot 8 \cdot 7}\right) = -\frac{1}{\sqrt{2\pi}}\cdot  \frac{479}{560} \approx -0.34124.\]
    \end{red}
    \item Find $P(b) - P(a)$. (Your answer should be close to 68\%.)
    
    \begin{red}
    \[P(1) - P(-1) \approx \frac{1}{\sqrt{2\pi}}\cdot  \frac{479}{560} - \left( - \frac{1}{\sqrt{2\pi}}\cdot  \frac{479}{560}\right) \approx 0.34124 + 0.34124 = 0.68248.\]
    \end{red}
    \item How many terms would be necessary to ensure that your approximation to $P(b)-P(a)$ is accurate to four decimal places? 
    \begin{red}
    \[
\begin{array}{c|c|c|c|c|c|c}
k     & 1         & 2         & 3         & 4         & 5         & 6         \\\hline
|a_k| & 0.1329808 & 0.0199471 & 0.0023747 & 0.0002309 & 0.0000189 & 0.0000013
\end{array}
    \]
    So the first term that's smaller than $5\times10^{-5}$ is $|a_5|$, which means we need $k=4$, which means we need five terms. (Remember, $k$ starts counting at 0.)
    \end{red}
\end{enumerate}
\end{document}